\documentclass[11pt, letterpaper]{article} % Se define el papel aquí para evitar conflictos

% --- CONFIGURACIÓN DE FUENTES ---
\usepackage[T1]{fontenc}
\usepackage{lmodern}
\usepackage{float}
\pdfinclusioncopyfonts=1

% --- PAQUETES DE FORMATO ---
\usepackage[margin=1in]{geometry} % Cargado SOLO UNA VEZ
\linespread{1.2}

% Un poco de espacio entre párrafos
\setlength{\parskip}{0.5\baselineskip}
\setlength{\parindent}{0pt}

% Math (AMS) DEBE IR ANTES DE newtxmath para evitar choque \Bbbk
\usepackage{amsmath,amssymb}

% Fuente matemática/texto
\usepackage{newtxtext,newtxmath}

\usepackage{graphicx}
\usepackage{xcolor} % required for \textcolor[HTML]{...} in figure legends
\definecolor{docnotelinkcolor}{rgb}{0,0,0} % link color used by naturemag.bst (DOI/URL); kept black to match hyperref scheme
\frenchspacing

% Tables
\usepackage{booktabs}
\usepackage{threeparttable}
\usepackage{float}
\usepackage{placeins}
\usepackage{caption}

% Math extra
\usepackage{bbm}
\setcounter{secnumdepth}{4}

% Citations
\usepackage[numbers,sort&compress]{natbib}
\usepackage{url}

\usepackage{hyperref}
\hypersetup{colorlinks=true,linkcolor=black,citecolor=black,urlcolor=black}
\captionsetup{font=scriptsize, labelfont=bf}

% ---- Macros ----
\newcommand{\parencite}[1]{\cite{#1}}
\newcommand{\textcite}[1]{\cite{#1}}

\title{\textbf{A Matthew Effect in Occupational Skill Content: Divergent Diffusion of Cognitive and Physical Skills Reinforces Inequality}}
%\author{
%Roberto Cantillan$^{1}$ \quad Mauricio Bucca$^{1,*}$\\[3pt]
%\small $^{1}$Department of Sociology, Pontificia Universidad Católica de Chile, Santiago, Chile\\
%\small $^{*}$Corresponding author. Email: mebucca@gmail.com
%}
\date{}

\begin{document}
\maketitle


\begin{abstract}
Technological change is expected to spread cognitive skills throughout the labor market, particularly into the lower-status occupations most exposed to automation and under the greatest pressure to upskill. We show the opposite. Cognitive skills systematically diffuse upward through the occupational hierarchy, whereas physical skills diffuse downward. Analyzing approximately 40 million directed diffusion opportunities spanning 741 occupations and 160 skill requirements in O*NET (2015--2024) using a directional gravity model, we isolate the effect of occupational status differences net of occupation-specific propensities to emit and receive skills and skill-specific diffusibility. Socio-cognitive skills are preferentially adopted by occupations above their current holders and abandoned by those below, producing a net upward diffusion. Sensory-physical skills exhibit the reverse pattern, moving downward through preferential adoption by lower-status occupations and abandonment by higher-status ones. Thus, rather than diffusing broadly across the occupational structure, skill diffusion follows divergent directions across skill types that reproduce occupational inequality by concentrating socio-cognitive skills in higher-status occupations and sensory-physical skills in lower-status ones---a Matthew effect operating among occupations themselves.
\end{abstract}

\clearpage

\section*{Introduction}

Contemporary labor markets are sharply polarized along skill-defined lines. One segment concentrates specialized socio-cognitive skills---analytical reasoning, complex problem-solving---alongside high educational requirements and high wages; the other relies on sensory-physical skills---manual dexterity, physical coordination---requires fewer credentials, and offers lower pay \parencite{liu_payoff_2013, anderson_skill_2017, alabdulkareem_unpacking_2018, hosseinioun_skill_2025}. Its broad contours align with a familiar economic account: technological change has raised the value of non-routine cognitive work while eroding the value of routine tasks---cognitive and manual alike---that can be codified and automated, growing employment and wages at the cognitive top, expanding manual and service work that resists automation, and hollowing out the middle \parencite{autor_skill_2003, acemoglu_skills_2011, acemoglu_robots_2020, zhang_future_2024}.

Technology is not only reshaping who works where; it is redefining the content of work itself---the bundles of capabilities and tasks that constitute an occupation \parencite{weeden_case_2005, labussiere_are_2026}. New technology, and most recently generative AI, may push all occupations to rely more heavily on cognitive skills and non-routine tasks \parencite{autor_skill_2003, spitz_oener_technical_2006, oesch_upgrading_2011, eloundou_gpts_2024, noy_experimental_2023}. This pressure is expected to fall hardest on lower-status occupations, which are most exposed to automation and most incentivized to emulate higher-productivity practices \parencite{acemoglu_robots_2020, goos_routine_2014, autor_skill_2015}; evidence from job postings supports this picture, with upskilling and skill diversification sharpest at the bottom of the wage distribution \parencite{tong_lower-skilled_2025, han_skill_2026}.

We show that the opposite is happening. Cognitivization is not spreading evenly, still less from the bottom up: socio-cognitive skills are adopted by occupations ranking \emph{above} their current holders and shed by those ranking below---flowing upward---while sensory-physical skills move the other way, adopted below their holders and shed from above. Because both flows run in the same direction rather than offsetting, the same upskilling pressure pulls the two ends of the occupational hierarchy apart rather than together, reproducing the cognitive--physical divide through a dynamic of cumulative advantage analogous to the Matthew effect \parencite{Merton1968}. For debates on AI, automation, and the future of work, the consequential question is therefore not whether work is changing but \emph{how and for whom}.

\subsection*{Theoretical framework}

Pressure to upskill is not the same as realized change in skill content. Whether, how, and for whom skill change occurs depends on an occupation's position within a relational domain \parencite{cheng_flows_2020, lin_network_2022, rytina_network_2020}, in which task similarity, training proximity, and productive complementarities---together with processes of emulation and differentiation---facilitate the diffusion of certain skills through specific channels while blocking others \parencite{strang_institutional_1993, hidalgo_product_2007, wimmer_domains_2021, bail_prestige_2019}, much as ideas and practices spread through networked social systems \parencite{watts_cascades_2002, banerjee_diffusion_2013, vosoughi_spread_2018}.

We study skill diffusion as a form of assortative mixing \parencite{newman_assortative_2002}---the tendency of network nodes to connect preferentially with similar others---but one that is asymmetric with respect to status. Skills spread toward the occupations that offer the strongest structural fit with their current holders---defined by symmetric features such as task similarity, productive complementarity, and training overlap \parencite{hidalgo_product_2007}. But fit alone does not explain the direction of skill spread. A nurse is more likely than a distant occupation to take on a skill established among physicians---linked by overlapping tasks, similar training, and routine contact---yet a nurse's propensity to adopt or shed a skill characteristic of physicians is not the mirror image of a physician's propensity to adopt or shed one characteristic of nurses. Diffusion depends, in other words, on the signed status distance between an occupation and its neighbors \parencite{bruch_aspirational_2018}: the same relatedness, we argue, yields different diffusion probabilities depending on which occupation ranks higher. We therefore study not just whether skills spread, but \emph{which skills, and in which direction, they travel} relative to the status of the occupations that already hold them.

\begin{figure*}[tp]

\centering

\includegraphics[width=.88\textwidth]{figures/illustration.png}

\caption{\textbf{The two elementary events the analysis is built from.} In an \emph{adoption} event, a target occupation $j$ that lacks a skill comes to hold it under positional pressure from source occupations $i$ that already do; in an \emph{abandonment} event, $j$ sheds a skill that source occupations $i$ retain. The analysis aggregates millions of such directed dyads to ask, for each event, not merely whether a skill moves but on which side of the status gap it lands. We do not observe transmission along particular dyads; we observe that an occupation's likelihood of gaining or losing a requirement moves with the status configuration of all the occupations that hold it.}

\label{fig:diffusion-schematic}

\end{figure*}


We formalize this process with a directional gravity framework that extends established models of trade, migration, and knowledge flows \parencite{tinbergen_shaping_1962, hidalgo_product_2007} to skill diffusion, capturing the aggregate positional pressure that the field of neighboring occupations collectively exerts on a skill profile---a diffuse influence with no single traceable source. The probability that a skill is adopted rises with the mass of both source and target occupations---their capacities to emit and absorb skills---and decays with the effective distance between them; abandonment likewise depends on both masses but rises with distance (Fig.~\ref{fig:diffusion-schematic}). Our extension makes this distance \emph{asymmetric}, adding the signed status difference between occupations to symmetric task-profile relatedness (see Methods). A fixed-effects strategy isolates the status-gap channel from occupations' compositional characteristics and skills' idiosyncratic features.

The model captures the joint signature of several relational mechanisms driving skill diffusion. One, \emph{emulation}, runs from the bottom up: occupations may take on the work content of higher-status neighbors and shed content tied to those below. But not all occupations are equally positioned to emulate, and several further mechanisms run the opposite way, channeling cognitive content toward the top and physical content down. \emph{Task allocation}: when firms route complex, non-routine work by reference to where it already sits, tasks bundled with high-status occupations are retained and deepened there, while the same tasks, once they surface in low-status work, are treated as offloadable \parencite{autor_skill_2003}. \emph{Absorptive capacity}: a requirement held by high-status occupations comes with the training, vocabulary, and discretion that make it legible to similarly positioned occupations; carried by low-status work, it lacks that scaffolding and is harder to retain \parencite{cohen_absorptive_1990}. \emph{Credentialist closure}: held above, a requirement is enclosed and protected by licensing and credentials; held below, it commands no such protection and is more readily shed \parencite{abbott_system_1988, weeden_why_2002, kleiner_licensing_2013}. \emph{Valuation asymmetry}: a cognitive requirement carried by high-status occupations signals competence and commands returns, inviting emulation, while the same requirement in low-status work is devalued---not because its content changed, but because of who holds it \parencite{deming_growing_2017}. Taken together, these mechanisms would leave lower-status occupations ill-equipped to absorb new cognitive skills while structurally priming higher-status ones to integrate them.

To trace this process we draw on O*NET, the U.S.\ Department of Labor database that records, for each occupation, the skills and capabilities its work requires. Our data span roughly 40 million directed diffusion opportunities---about 21.5 million adoption and 18.6 million abandonment dyads---across 741 occupations and 160 skill requirements between 2015 and 2024. Comparing skill profiles over time, we observe which occupations take on new skills and which let old ones lapse, and where each skill's current holders sit in the status hierarchy relative to the occupation that adopts or sheds it. We measure status as a composite of occupational standing---combining education, earnings, and cognitive-skill reliance, all fixed at their 2015 baseline---so that the directional results below describe which occupations subsequently gain or shed skills, not how status itself was defined (see Methods).


\section*{Results}

We analyze skill adoption and skill abandonment across three functional types defined by established taxonomies. We distinguish by functional domain (socio-cognitive vs.\ sensory-physical) \parencite{alabdulkareem_unpacking_2018} and functional specificity (general vs.\ specialized) \parencite{hosseinioun_skill_2025}---which empirically yields three skill types: general socio-cognitive (Gen.\ SC), specialized socio-cognitive (Spec.\ SC), and sensory-physical (specialized sensory-physical is empirically absent; see Methods). This taxonomy allows us to test whether diffusion patterns are uniform across skills or vary systematically by type and domain. Throughout, an occupation \emph{adopts} (or \emph{abandons}) a skill when its requirement profile crosses (or drops below) the threshold of revealed specialization in that skill between 2015 and 2024; \emph{diffusion} denotes this directional updating relative to occupations that already hold the skill, not observed transmission between specific occupational pairs (see Methods).

\textbf{Descriptive results: cognitive skills flow upward, physical skills flow downward} Change in occupational skills is structured by two forces: one symmetric (relatedness) and one asymmetric (status gaps). The diffusion pattern, however, diverges across cognitive and physical functional domains, in both adoption and abandonment flow. On the symmetric side, occupations are more likely to incorporate skills already present in occupations with similar skill profiles, and less likely to borrow from dissimilar ones---a pattern that holds for physical and cognitive skills alike, both general and specialized. Abandonment mirrors this logic: occupations are more likely to shed skills still held by occupations with dissimilar profiles, and to retain those held by similar ones (Fig.~\ref{fig:flows_gradients}A--B).

Yet relatedness captures only one dimension of the process. Whether an occupation sits above or below the potential source of a skill also shapes how skills move—and it does so in opposite directions for socio-cognitive and sensory-physical skills. In adoption, socio-cognitive skills are disproportionately taken up by occupations that rank above those already holding them, flowing upward through the hierarchy---and more strongly for specialized than for general socio-cognitive skills. Sensory-physical skills show the mirror image: their adoption grows more likely the further the adopter sits below the occupation previously holding the skill, so they are taken up disproportionately below their current holders and move predominantly downward (Fig.~\ref{fig:flows_gradients}C).

Abandonment provides a complementary picture for assessing whether skill loss offsets this divergence or amplifies it. A corrective process would erode the adoption-driven pattern—shedding cognitive skills from high-status occupations, or retaining them in low-status ones. Instead, abandonment compounds the divergence rather than offsetting it. Socio-cognitive skills are more likely to be shed when the occupation dropping them ranks below the occupation that still holds them—a pattern observed for both specialized and general skills. Sensory-physical skills show the opposite: they are disproportionately dropped by occupations that rank above those still holding them, so physical content tends to persist in lower-status profiles (Fig.~\ref{fig:flows_gradients}D).

These descriptive gradients establish the empirical puzzle. Skill diffusion is not a generalized process of cognitive upgrading but one that reproduces the cognitive--physical divide, channeling socio-cognitive skills toward higher-status occupations and sensory-physical skills toward lower-status ones.

\begin{figure*}[tp]
\centering
\includegraphics[width=.88\textwidth]{figures/fig1_full_networks.pdf}
\caption{%
\textbf{Skill adoption and abandonment rates as a function of skill-profile distance and occupational status gap.}
Colors distinguish three skill types throughout: \textcolor[HTML]{3B4992}{\textbf{general socio-cognitive}}, \textcolor[HTML]{008280}{\textbf{specialized socio-cognitive}}, and \textcolor[HTML]{EE0000}{\textbf{sensory-physical}}. In all panels, solid dots represent binned mean rates across 25 equal-frequency bins, shaded ribbons show 95\% bootstrap confidence intervals, and solid curves display LOESS-smoothed trends.
\textbf{(A, B)}~Adoption (A) and abandonment (B) rates plotted against the cosine distance between the skill-specialization profiles of the source and target occupations (higher values = less related).
\textbf{(C, D)}~Adoption (C) and abandonment (D) rates plotted against the composite PC1 status gap (target minus source; see Methods). Negative values indicate the target ranks below the source (downward gap), while positive values indicate the target ranks above (upward gap). The dashed vertical line marks zero.
\textbf{Insets:} The network diagrams in panels C and D show directional flow rates aggregated to status quintiles (Q1 = lowest, Q5 = highest). Each arc connects a source quintile to a destination quintile, with arc width and opacity proportional to the normalized directional flow rate within each skill class. Darker, right-pointing arcs indicate upward flows; lighter, left-pointing arcs indicate downward flows. Subtitle percentages report the overall mean upward ($\uparrow$) and downward ($\downarrow$) cross-quintile rates computed over the respective risk sets. O*NET, 2015--2024.
}
\label{fig:flows_gradients}
\end{figure*}


\textbf{Fixed-effects gravity models: status gaps drive diffusion, net of occupational and skill heterogeneity.} The observed divergence need not reflect status differences. It could instead arise from stable features of occupations and skills: some occupations act more as sources of influence than others (for instance, large, central, or early-adopting ones); some adopt or abandon skills more readily overall (for instance, those reorganizing rapidly or facing volatile task demands); and some skills simply travel more easily than others (for instance, general-purpose rather than narrowly specialized ones). Each may act as a confounder: if high-status occupations are voracious adopters in general, their apparent pull on cognitive skills could be an artifact of that appetite rather than evidence of directional, status-based diffusion. To separate the status gap effect from such alternatives, we estimate gravity models for both adoption and abandonment under two complementary fixed-effect strategies, each identifying the status effect from a different source of variation (see Methods). The target-varying specification (source fixed effects; Panels~A and C) absorbs all stable characteristics of the source occupation and the skill, identifying the status effect from variation in the target's position relative to a fixed source. The source-varying specification (target fixed effects; Panels~B and D) instead absorbs all stable characteristics of the target occupation and the skill, identifying the effect from variation in the source's position relative to a fixed target; because this removes the target's absolute status and its overall propensity to adopt or shed, it shows the effect reflects the signed status gap rather than the focal occupation's general activity level. Both specifications further condition on skill-profile relatedness, so that status effects are net of ordinary proximity between occupational profiles—ensuring that directional friction indexes movement along the status hierarchy, not profile similarity. 

\begin{figure*}[tp]
\centering
\includegraphics[width=0.90\linewidth]{figures/fig2_regression.pdf}
\caption{%
\textbf{Fitted log-hazard of skill adoption and abandonment as a function of signed status gap, under two complementary fixed-effect strategies.}
Colors distinguish three skill types throughout: \textcolor[HTML]{3B4992}{\textbf{general socio-cognitive}}, \textcolor[HTML]{008280}{\textbf{specialized socio-cognitive}}, and \textcolor[HTML]{EE0000}{\textbf{sensory-physical}}. The horizontal axis across all panels represents the composite PC1 status gap (target minus source). The unshaded region (left of $G_{ij}=0$) corresponds to downward gaps (target ranks below source), while the gray-shaded region (right) corresponds to upward gaps (target ranks above source). The vertical axis is the log-hazard from the complementary log-log gravity model (Eq.~\ref{eq:gravity_main}).
\textbf{(A, B)}~Skill adoption under source and skill fixed effects (A), identifying the gap effect from variation across targets, and under target and skill fixed effects (B), identifying it from variation across sources.
\textbf{(C, D)}~Skill abandonment under the same two strategies: source and skill fixed effects (C) and target and skill fixed effects (D).
\textbf{Visual encodings:} Solid lines represent the fitted piecewise linear predictor evaluated along the status gap. Open circles mark predicted values at 2.5-unit reference intervals ($-5$, $-2.5$, $0$, $+2.5$, $+5$). In-plot text annotations report the estimated slopes for downward ($\hat{\beta}^{\downarrow}$) and upward ($\hat{\beta}^{\uparrow}$) gaps by skill class, with standard errors in parentheses (clustered three-way by source occupation, target occupation, and skill). The rug plot beneath the horizontal axis displays the distribution of dyad-level status gaps for a random subsample of 60{,}000 dyads per flow. Adoption: $n\approx21.5$\,M directed dyads; Abandonment: $n\approx18.6$\,M. O*NET, 2015--2024.
}
\label{fig:model_projection_combined}
\end{figure*}

 
The adoption estimates reveal a clear directional signature (Fig.~\ref{fig:model_projection_combined}, top row). In the target-varying specification (source FE; Panel~A), socio-cognitive skills are taken up most readily when the adopting occupation ranks above the occupation that already holds the skill—when the skill travels up the status gap—and this upward pull is strongest for specialized cognitive skills. Sensory-physical skills show the mirror pattern: they are taken up when the adopter ranks below the prior holder, and face a steep barrier to upward movement. The source-varying specification (target FE; Panel~B) reproduces this configuration, confirming that the gradient tracks the direction of the gap itself, not the prior holder's general tendency to emit skills, the adopter's appetite for new ones, nor---given the skill fixed effects---the intrinsic diffusibility of the skills. Adoption thus runs up the gap for cognitive skills and down it for physical ones. The two domains also differ in how the upward and downward slopes compare. For socio-cognitive skills both slopes are positive under source fixed effects---adoption rises with the target's status whether the target sits above or below the current holder---though under target fixed effects the downward branch flattens to near zero for general socio-cognitive skills. For sensory-physical skills the slopes differ sharply---a steep barrier to upward adoption against a near-flat downward branch---meaning that these skills are taken up mainly by occupations ranking at or below their current holders and seldom move upward.

Abandonment runs in the complementary direction along the same gap, compounding the divergence (Fig.~\ref{fig:model_projection_combined}, bottom row). In the target-varying specification (source FE; Panel~C), a sensory-physical requirement is most likely to be dropped when the occupation shedding it ranks above the occupation with which it shared the skill; this upward release of physical content is, if anything, steeper under the source-varying specification (target FE; Panel~D). Socio-cognitive skills run the other way: cognitive content is dropped down the gap—shed when the shedding occupation ranks below the occupation it shared the skill with—and retained when the gap runs upward. This cognitive gradient is more muted under the source-varying specification (Panel~D), indicating that some of the cognitive retention reflects durable features of the holding occupations alongside the status gap itself. In addition, the positive boundary-term $\kappa$ indicates that occupations are more likely to drop a physical skill discretely once they rank above the occupation they shared it with, rather than only as the gap widens.

Read together, the two flows trace a single relational pattern: cognitive skills move up the gap in adoption and down it in abandonment, while physical skills move down the gap in adoption and up it in abandonment. The two are not inverses of one another but two arms of one process, reproducing the cognitive–physical divide through what occupations add and what they let go alike. These estimates identify the direction of skill flow as a function of relative status, net of stable occupation- and skill-level heterogeneity and of task relatedness (the individual triads underlying these estimates can be examined via the interactive explorer in \hyperref[sec:SI_triads]{Supplementary Section~\ref*{sec:SI_triads}}).

These patterns are robust to how skill types are measured and to how status is defined. They do not stem from the relative (RCA) construction of the specialization index (\hyperref[sec:SI_rca_denom]{Supplementary Section~\ref*{sec:SI_rca_denom}}), are robust to moderate skill-type misclassification (\hyperref[sec:SI_missclass]{Supplementary Section~\ref*{sec:SI_missclass}}), and do not rest on any single dimension of the composite status measure---log wage, log education, and cognitive content each reproduce the same directional asymmetry on their own (\hyperref[sec:SI_alt_status]{Supplementary Section~\ref*{sec:SI_alt_status}}). In particular, wage and education alone---neither of which enters any skill outcome---reproduce the upward pull on cognitive skills, so the pattern is not an artifact of cognitive content appearing in both the status index and the outcome.

\vspace{10pt}

\textbf{Skill-level divergent diffusion reproduces aggregate occupational stratification.} If dyadic requirement diffusion is stratifying, its accumulated effects should recover the aggregate gradient of occupational requirement change. We test this implication by projecting the estimated directional channel parameters onto the full occupational risk set and comparing the resulting directional profile with two nested counterfactuals: a symmetric status null that retains the magnitude of the status gap but ignores its direction, and a distance-only null that ignores status altogether and retains only skill profile proximity. The projection is primarily a test of directional sufficiency rather than a calibrated fit: it asks whether the signed status gap alone, carried from the dyadic estimates to the full risk set, suffices to regenerate the observed gradient. It largely does---the directional profile recovers most of the socio-cognitive macro-gradient and about half of the sensory-physical one under source fixed effects---whereas both direction-blind nulls stay close to flat, confirming that neither proximity nor the size of the status gap---absent its direction---can account for the observed gradient. The directional rule documented at the dyadic level scales to the observed macro-pattern of occupational skill consolidation \parencite{blau_inequality_1977}: socio-cognitive skills become increasingly concentrated in higher-status occupations, while sensory-physical skills become concentrated in lower-status ones through adoption and abandonment operating together (Fig.~\ref{fig:projection_main}).

On adoption, both socio-cognitive skill types rise monotonically across status quintiles, while sensory-physical skills decline in mirror image (Fig.~\ref{fig:projection_main}, adoption row in both panels). The directional projection recovers the sign of every observed Q5--Q1 gradient under both fixed-effect strategies, and most of its magnitude for the socio-cognitive classes under source fixed effects (Panel~A)---about half for sensory-physical; under target fixed effects (Panel~B) the socio-cognitive magnitude is only partially recovered, as the gap coefficients attenuate once the target's own status is absorbed. Both nested counterfactuals remain comparatively flat throughout, confirming that the direction of the status gap, not its magnitude, accounts for the pattern (full statistics in Table~\ref{tab:SI_gradients}).

On abandonment, requirement loss reinforces the same hierarchy: both socio-cognitive classes show declining abandonment across status quintiles, while sensory-physical abandonment rises in mirror image as higher-status occupations disproportionately shed physical content (Fig.~\ref{fig:projection_main}, abandonment row in both panels). The directional model recovers the sign of these gradients under both strategies and most of their magnitude for the socio-cognitive classes under source fixed effects (Panel~A), with the socio-cognitive magnitude only partially recovered under target fixed effects (Panel~B); the sensory-physical gradient, by contrast, is recovered about as fully under target as under source fixed effects. The symmetric status null recovers only a small fraction of the gradient for either socio-cognitive class---confirming that direction, not merely the size of the status gap, drives the stratifying pattern.

\begin{figure*}[tp]
\centering
\includegraphics[width=0.85\textwidth]{figures/fig3_unified.pdf}
\caption{%
\textbf{The dyadic directional rule reproduces the aggregate gradient of occupational skill change.}
Observed and model-projected adoption (left) and abandonment (right) rates by occupational status quintile (Q1 = lowest, Q5 = highest) for each skill class, under both complementary fixed-effect strategies: source and skill fixed effects (\textbf{Panel A}, top) and target and skill fixed effects (\textbf{Panel B}, bottom). The directional model applies the full set of estimated channel parameters to every dyad in the risk set and averages within quintiles; the symmetric null retains only the magnitude of the status gap, and the distance-only null only task proximity. Socio-cognitive skills concentrate in higher-status occupations and sensory-physical skills in lower-status ones, and the directional model tracks this gradient under both strategies while the direction-blind nulls stay close to flat. Full Q5--Q1 gradient statistics for both strategies are reported in Table~\ref{tab:SI_gradients} (\hyperref[sec:SI_projections]{Supplementary Section~\ref*{sec:SI_projections}}). O*NET, 2015--2024.
}
\label{fig:projection_main}
\end{figure*}


 
\section*{Discussion}

The standard upgrading account predicts that technological pressure should extend cognitive skill requirements across the labor market, including downward to the occupations most exposed to automation and most pressed to upskill. We find the opposite. Realized skill change does not diffuse cognitive content downward; it sorts it by status. Socio-cognitive skill requirements are taken up preferentially by occupations ranking above their current holders and shed by those ranking below, while sensory-physical skills move in the reverse direction—adopted downward and shed upward. Adoption and abandonment compound rather than offset, each pushing cognitive content up the status gap and physical content down it. The same upskilling pressure thus tracks divergence rather than convergence, as a function of the signed status distance between occupational pairs. And although that pressure appears sharpest at the bottom \parencite{tong_lower-skilled_2025, han_skill_2026}, cognitive skills accrue at the top nonetheless---suggesting that what propagates across occupational profiles depends less on where pressure is greatest than on where status already sits.

Both task proximity and status shape diffusion, but only the signed status gap can account for the cognitive--physical divergence. Knowing how far apart two occupations stand, and how large the status gap between them is, but not which one stands higher, is not enough to reproduce the divergent structure of skill change; knowing the direction is. The evidence bears this out: specialized socio-cognitive content flows upward far more strongly than its sheer availability would predict, while socio-cognitive abandonment follows the direction of rank rather than the size of the gap---which, under source fixed effects, accounts for almost none of it. The directional signature holds under both fixed-effect strategies, even where its estimated magnitude attenuates. This directional pull is not an artifact of which occupations happen to change most or which skills travel most easily---those general tendencies are held constant---and it is difficult to attribute to a technological shock shared by neighboring occupations, since such a shock would tend to push in one direction across the board rather than reverse itself across skill domains. These results identify the direction of skill flow as a function of relative status; they are not estimates of a causal effect of status itself, nor do they adjudicate among the relational mechanisms---emulation, task allocation, absorptive capacity, credentialist closure, valuation asymmetry---that plausibly produce it. What they pin down is a specific structure---the direction of skill flow along the status gap, its reversal across domains, and its persistence net of stable occupation- and skill-level propensities---none of which follows from the static distribution of skills or is readily reproduced by a correlated shock; which of these mechanisms supplies the motive force is a separate question, one for data at the organizational or individual level.

Occupational inequality, on this account, is reproduced not only because occupations differ in who works in them, how much those workers earn, and what the work requires, but because the skills occupations gain and lose travel in systematically different directions. This suggests how the skill polarization documented in prior work \parencite{alabdulkareem_unpacking_2018} is sustained: rather than a fixed feature of the occupational structure, polarization is continually renewed as occupations that already hold cognitive skills draw still more of them while shedding physical content, and occupations built on physical skills retain it while cognitive content departs. The pattern we estimate is a cumulative-advantage dynamic in skill space---a Matthew effect written into the architecture of occupations themselves.

These results also help reconcile a tension in recent evidence on skill change. Within-job skill diversification is fastest at the bottom of the wage distribution, yet the returns to it are lowest precisely there \parencite{han_skill_2026}. Our findings suggest why: the skills that actually accrue to lower-status occupations are largely confined to the sensory-physical domain, where diffusion concentrates regardless of organizational intent.

A second tension concerns upskilling pressure itself. Evidence that lower-skilled occupations face the steepest such pressure \parencite{tong_lower-skilled_2025} captures aspirational upgrading---what organizations signal they want. What we observe is which skills actually propagate across occupational profiles and persist in them, and the two diverge exactly where pressure is greatest: at the bottom of the status distribution, sensory-physical skills nonetheless move predominantly downward. Stated ambition, in other words, does not by itself determine what occupations become; the direction of skill flow is constrained by relational mechanisms our design isolates but cannot name. A plausible aggregate consequence---one our estimates point to rather than measure directly---is the qualitative depletion of middle-wage destinations: occupations that historically bundle cognitive monitoring with hands-on execution, such as routine coordination, clerical, and technical-support roles, are positioned to lose socio-cognitive skills upward to professional and managerial destinations while remaining receptive to sensory-physical skills shed from above. Operating through the composition of skills rather than the volume of workers, this process is distinct from routine-task displacement \parencite{autor_growth_2013, acemoglu_skills_2011}, and it qualifies the secular tendency for occupations to accrete tasks over time \parencite{jackson_division_2025}: that accretion is status-sorted rather than domain-neutral, so aggregate broadening can coexist with systematic divergence across the hierarchy.

Theoretically, the pattern extends relational accounts of stratification \parencite{tomaskovic-devey_relational_2019, tahlin_skill_2023, cheng_flows_2020, lin_network_2022} from the sorting of workers into positions to the sorting of skills across the positions themselves. It is consistent with the gradationalism documented in occupational mobility \parencite{york_gradationalism_2025}, yet reveals an asymmetry a continuous mobility metric cannot recover: signed status distance produces direction-specific decay, a categorical-like barrier within a formally continuous space. Where prior work locates the reproduction of status boundaries in hiring networks and in how occupations expand into adjacent domains \parencite{clauset_systematic_2015, henaut_polyoccupationalism_2023}, our results locate a related boundary earlier still---in which skills an occupation gains and loses, before any worker is selected, hired, or paid.

Our results are relevant under accelerating AI adoption. Large language models expose roughly half the U.S.\ workforce to task-level substitution \parencite{eloundou_gpts_2024}, productivity gains from AI assistance concentrate in cognitively intensive tasks \parencite{noy_experimental_2023}, and substitution potential is highest for well-specified physical and routine cognitive tasks \parencite{acemoglu_robots_2020, zhang_future_2024}---precisely the skills we find most confined to the bottom of the hierarchy. Our results therefore imply a compounded barrier for workers displaced from physical-intensive occupations: the skills they hold are both most exposed to substitution and least able to cross the domain boundary separating them from cognitively intensive destinations. External evidence on worker mobility is consistent with such a blockage---transitions from declining to growing occupations account for only about 5\% of moves over two decades, with the odds of such a move less than a third of those of remaining in a declining sector \parencite{song_trapped_2026}. As AI augments productivity unevenly across the hierarchy and the returns to complementary cognitive skills concentrate at the top \parencite{acemoglu_kong_2026}, these trends seem more likely to reinforce the directional frictions we document than to dissolve them.

A few important limitations bound these conclusions. Our analysis characterizes skill diffusion at the occupational level; it does not observe the organizational decisions that produce individual adoptions, nor the general-equilibrium dynamics that accompany shifting skill demand \parencite{alabdulkareem_unpacking_2018}, and the depletion of middle-wage destinations is an interpretation of the estimated directional channel rather than a directly measured flow. The observational design cannot decisively rule out that some profile change reflects occupations responding independently to common technological pressures, though the domain-reversing structure of the gradient---cognitive skills rising precisely where physical skills fall---sharply constrains that reading. Finally, the analysis rests on U.S. data, and the magnitude of directional channeling is likely to be institutionally contingent: directional barriers should weaken where credentialing is fungible, occupational boundaries are permeable, and complementary skills are easier to bundle across domains. Cross-national tests exploiting variation in credentialing regimes, wage-setting structures, and union density are needed to separate general patterns from institutional contingencies \parencite{hosseinioun_skill_2025}. That contingency is also where leverage lies: interventions that build credential bridges or stabilize cross-domain skill complements target the directional friction itself, a mechanism our estimates suggest upskilling pressure alone does not reach \parencite{moro_universal_2021}.

Occupations, in this light, are not fixed bundles of tasks but moving targets whose composition is sorted by status before any worker enters them. Worker mobility traces the boundaries that structure careers; skill diffusion traces the boundaries that structure what occupations become.

\section*{Methods}

\subsection*{Data and measures}
 
We draw occupation-level skill requirement data from the Occupational Information Network (O*NET) across two extracts anchored at baseline ($t_0 = 2015$) and endline ($t_1 = 2024$). O*NET reports the importance and required level of 160 skill, knowledge, and ability items spanning six content domains for 873 standardized occupations via a rolling update design in which approximately one-fifth of occupations is resurveyed each year \parencite{alabdulkareem_unpacking_2018}. We restrict our analysis to the Skills, Knowledge, and Abilities domains, which together constitute the structural capacity that occupations signal they require. After taxonomic harmonization and restriction to occupations with valid profiles at both endpoints---procedures detailed in \hyperref[sec:SI_onet]{Supplementary Section~\ref*{sec:SI_onet}}---the analytic sample covers 741 occupations (see \hyperref[sec:SI_sample]{Supplementary Section~\ref*{sec:SI_sample}} for the full derivation of the analytic sample). We use an endpoint window rather than annual panels to match this structure: because only a fraction of occupations is resurveyed in any given year, intermediate-wave non-change reflects measurement timing as much as substantive stability. The nine-year window is long enough that all 741 occupations receive at least one full survey refresh within the horizon, allowing both adoption and abandonment to be measured cleanly at endline without imputing annual timing. Throughout, \textit{diffusion} denotes directional updating of occupational skill requirements under a risk-set design with a plausible source in adoption and a plausible positional referent in abandonment; we make no claim about direct transmission events between specific occupational pairs.
 

\textbf{Wages, education, task-based socio-cognitive intensities}: Occupational wages are drawn from the Bureau of Labor Statistics Occupational Employment and Wage Statistics (OEWS), which reports median annual wages for each of the 873 O*NET occupations. Educational requirements are operationalized as the occupation-level expected level of required education on the O*NET twelve-category \emph{Required Level of Education} scale (Scale~RL). Cognitive task content is the share of an occupation's total O*NET skill-importance accounted for by the socio-cognitive skill items identified in the domain classification described below. All three variables are constructed from 2015 values and held fixed throughout the observation window---predetermined with respect to outcomes---to preclude reverse causality from skill adoption mechanically attracting higher wages or educational requirements. Full variable definitions and descriptive statistics are reported in \hyperref[sec:SI_wages]{Supplementary Section~\ref*{sec:SI_wages}}.

\textbf{Skill classification by functional domain and specificity.}: We classify each of the 160 skill requirements along two orthogonal dimensions: functional domain and functional specificity. Functional domain distinguishes socio-cognitive from sensory-physical skills based on the community structure of the skill complementarity network, following \textcite{alabdulkareem_unpacking_2018}. For each occupation--skill pair we compute a revealed comparative advantage (RCA) score that captures whether a given skill is more intensively required by that occupation than would be expected from the economy-wide distribution of skill demands; we threshold at RCA~$\geq 1$ to identify specialization and apply Louvain community detection to the resulting co-specialization network (we obtain identical results using the Leiden algorithm and nearly identical results using the walktrap algorithm; results available upon request). The algorithm recovers a stable two-community solution in which one cluster concentrates social, cognitive, and knowledge-intensive items and the other concentrates physical, sensory, and manual items; we assign domain labels accordingly. Critically, the same two-community solution obtains when we apply the procedure to the 2024 endline data; we use the 2015 classification throughout to ensure that domain assignment is predetermined with respect to outcomes (\hyperref[sec:SI_domain]{Supplementary Section~\ref*{sec:SI_domain}}). Functional specificity captures each skill's structural position within the occupation--skill matrix. We characterize each skill $s$ by a nestedness contribution score $c_s$, defined as the standardized change in a matrix-level nestedness index when skill $s$ is removed from the baseline occupation--skill matrix \parencite{hosseinioun_skill_2025}. Intuitively, a skill with high $c_s$ is broadly required across many occupations in a hierarchically ordered pattern and scaffolds the skill hierarchy; a skill with low $c_s$ is more narrowly concentrated and depends on that broader foundation. Within the socio-cognitive domain, all skills have $c_s > 0$ by construction---the Louvain community that defines the cognitive cluster captures the nested core of the skill space---so we split cognitive skills at the within-domain median of $c_s$: those below the median are classified as \textit{specialized socio-cognitive} skills (narrowly concentrated cognitive skills) and those at or above as \textit{general socio-cognitive} skill requirements (broadly scaffolding cognitive skills). Sensory-physical skills are classified as \textit{sensory-physical} across all nestedness levels; the high-nestedness sensory-physical cell is empirically absent, and results are unchanged when it is retained as a separate category (\hyperref[sec:SI_taxonomy]{Supplementary Section~\ref*{sec:SI_taxonomy}}). Both $c_s$ and domain classification are constructed from 2015 values and are therefore predetermined with respect to outcomes. Full derivation of RCA, $c_s$, and the complete three-class taxonomy are reported in \hyperref[sec:SI_skill_class]{Supplementary Sections~S1.2.1--S1.2.4}.
 
\textbf{Occupational status index and status gap}: We summarize the status of each occupation in a scalar $\sigma_i$ and define the signed status gap as $G_{ij} = \sigma_j - \sigma_i$, where positive values indicate that the target occupation ranks higher than the source. We operationalize $\sigma_i$ as the first principal component of three occupation-level variables---median wage (log-transformed), the expected level of required education (log-transformed), and cognitive task content---estimated via PCA with mean-centering and unit-variance scaling across all 741 occupations in the analytic sample. PC1 accounts for the large majority of the joint variance across the three dimensions, with loadings of comparable magnitude across all three variables, confirming that they converge on a single underlying status axis. The signed status gap $G_{ij}$ is then computed as the difference in occupation-level scores: $G_{ij} = \sigma_j - \sigma_i$. All three components are measured at baseline (2015) and are therefore predetermined with respect to the adoption and abandonment outcomes, which are realized only by 2024. This temporal ordering is what keeps the status index from circularity: an occupation ranks high in $\sigma_i$ partly because it \emph{already} relied on cognitive content in 2015, so the result that cognitive skills travel toward higher-status occupations is a statement about which occupations subsequently \emph{gain} such content---not a definitional consequence of their baseline composition. Component variables, scaling decisions, and the full distribution of status scores are reported in \hyperref[sec:SI_pca]{Supplementary Section~\ref*{sec:SI_pca}}.

%more conceptual 
\textbf{Symmetric skill profile distance.}: Task relatedness between each source--target pair is measured as one minus the cosine similarity of their O*NET skill-specialization vectors at baseline, where each vector encodes binary RCA~$\geq 1$ indicators across the 160 skill items---that is, whether each occupation is more specialized in a given skill than the economy-wide average. Profile distance thus captures how similar two occupations are in what they require, independently of where they sit in the status hierarchy. Including it in the gravity model ensures that directional gap terms index movement along the status hierarchy rather than proxying for task similarity between occupational profiles. All distances are computed from 2015 skill vectors and are therefore predetermined with respect to outcomes.
 
\subsection*{Fixed-effect gravity models}
 
The unit of analysis is a directed triple $(i, j, s)$: source occupation $i$, target occupation $j$, and skill requirement $s$. We construct separate risk sets for the two diffusion flows, each exhaustive within its conditioning set to avoid selection on the dependent variable \parencite{bail_prestige_2019}. Because the risk set is defined relative to a source, a single change at the target---$j$ gaining or shedding $s$---enters as one triad for each occupation that held $s$ at baseline: a skill anchored in many occupations presents a denser field of diffusion opportunities than one held by few, and each source--target pairing is a distinct opportunity rather than a repeated record of the same event. The coefficients accordingly describe the structure of these opportunities---which configurations of relatedness and signed status are more likely to culminate in adoption or abandonment---and the three-way clustering by source, target, and skill accounts for the dependence this shared structure induces.
 
For adoption, a triple enters the risk set if $i$ is specialized in $s$ at baseline and $j$ is not, constituting a diffusion opportunity in which $s$ can plausibly propagate from $i$ to $j$. Applying this rule across all ordered occupation pairs and the full 160-skill taxonomy yields approximately 21.5 million adoption opportunities. Skill specialization is defined via revealed comparative advantage \parencite{alabdulkareem_unpacking_2018}: occupation $j$ is specialized in $s$ if $\mathrm{RCA}(j,s) \geq 1$, where RCA compares the relative importance of $s$ in $j$ to its economy-wide average, normalizing for skill ubiquity and identifying specialization relative to the aggregate distribution of skills rather than raw importance scores. The binary outcome $Y^{\mathrm{adopt}}_{ijs} = 1$ if target $j$ crosses the specialization threshold by 2024, conditional on not being specialized at 2015.
 
For abandonment, a triple enters the risk set if both $i$ and $j$ are specialized in $s$ at baseline. Given that source and target share a requirement at $t_0$, the question is whether target $j$ sheds it by $t_1$ as a function of the signed status gap $G_{ij} = \sigma_j - \sigma_i$. The outcome is located at the target: abandonment asks whether occupation $j$ drops a requirement it initially shared with $i$, with source $i$ entering only as a positional referent that locates $j$ in the status hierarchy. Applying this rule yields approximately 18.6 million abandonment opportunities. The binary outcome $Y^{\mathrm{aband}}_{ijs} = 1$ if target $j$ falls below the specialization threshold by 2024, conditional on being specialized at 2015.
 
We formalize the directional diffusion intuition through a gravity-hazard framework in which the adoption and abandonment hazards are jointly governed by occupational mass, skill-profile proximity, and signed status distance. For any skill $s$ and dyad $(i,j)$ in the flow-specific risk set $\Omega^{f}$, with $f \in \{\mathrm{adopt},\mathrm{aband}\}$, the cumulative hazard that a diffusion event occurs over the interval $(t_0,t_1]$ is governed by a common gravity kernel:
%
\begin{equation}
\Lambda^{f}_{ijs} \;\propto\; \frac{M_i \, M_j \, S_s}{D_{ij}}, \qquad (i,j,s)\in \Omega^{f}, \quad f \in \{\mathrm{adopt},\mathrm{aband}\}
\end{equation}
%
Here $\Lambda^{f}_{ijs}$ is the cumulative hazard of the event over $(t_0,t_1]$ within flow $f$; $M_i$ and $M_j$ capture the occupational mass of source and target, governing each node's capacity to influence and be influenced; $S_s$ captures skill-specific diffusibility; and $D_{ij}$ is an effective distance that combines the symmetric skill-profile distance $\mathrm{dist}_{ij}$ with the signed status gap $G_{ij}$, decomposed explicitly in Eq.~\ref{eq:gravity_main} below, rendering friction asymmetric across the occupational hierarchy. The adoption and abandonment risk sets $\Omega^{\mathrm{adopt}}$ and $\Omega^{\mathrm{aband}}$ are those defined above; in both, source $i$ acts as the positional referent and target $j$ is the focal occupation whose behavior is modeled, so the same kernel, evaluated over each risk set, generates the adoption and abandonment predictions within a unified framework.
 
We estimate this kernel with a complementary log-log (cloglog) link, whose linear predictor is exactly the log cumulative hazard, $\eta_{ijs} = \operatorname{cloglog}\big(P(Y_{ijs}=1)\big) = \log\big(-\log(1-P_{ijs})\big) = \log \Lambda^{f}_{ijs}$; over the single grouped interval $(t_0,t_1]$ this is the discrete-time proportional-hazards model \parencite{jenkins_easy_1995}. Taking logs of the kernel renders $\eta$ additively separable: the mass and diffusibility terms $\log M_i$, $\log M_j$, and $\log S_s$ are recovered by fixed effects $\alpha_i$, $\alpha_j$, and $\alpha_s$---analogous to multilateral resistance terms in structural trade gravity \parencite{fally_structural_2015, egger_glm_2015}---while the effective distance enters through $-\log D_{ij} = \beta^{\uparrow}_{g}\,\Delta^{\uparrow}_{ij} + \beta^{\downarrow}_{g}\,\Delta^{\downarrow}_{ij} + \kappa_{g}\,\mathbbm{1}[G_{ij}>0] + \delta_{g}\,\mathrm{dist}_{ij}$, which carries the friction parameters of interest. We estimate adoption and abandonment separately so that these parameters may vary by flow; in particular, $\delta_g$ is estimated separately by flow and takes opposite signs across them---negative for adoption, positive for abandonment---so the effective distance $D_{ij}$ is itself flow-specific, and the same profile dissimilarity that raises adoption friction lowers it for abandonment (adoption decays with distance while abandonment rises with it). The full linear predictor for triple $(i,j,s)$ is:
%
\begin{equation}
\eta_{ijs} = \underbrace{\alpha_i + \alpha_j + \alpha_s}_{\text{structural fixed effects}} + \underbrace{\beta^{\uparrow}_{g}\,\Delta^{\uparrow}_{ij} + \beta^{\downarrow}_{g}\,\Delta^{\downarrow}_{ij} + \kappa_{g}\,\mathbbm{1}[G_{ij}>0]}_{\text{asymmetric status friction}} + \underbrace{\delta_{g}\,\mathrm{dist}_{ij}}_{\text{skill profile distance}},
\label{eq:gravity_main}
\end{equation}
%
where $g \in \{\text{Spec.\ SC}, \text{Gen.\ SC}, \text{Phys}\}$ indexes the skill class, and $\Delta^{\uparrow}_{ij} = \max(0,\, G_{ij})$ and $\Delta^{\downarrow}_{ij} = \min(0,\, G_{ij})$ decompose the signed status gap into upward and downward components, allowing each direction of propagation to be governed by a distinct friction parameter without imposing symmetry. Estimating $\beta^{\uparrow}_{g}$ and $\beta^{\downarrow}_{g}$ separately lets upward emulation and downward offloading---potentially distinct processes---take different values rather than constraining them to mirror one another. The direction-entry indicator $\mathbbm{1}[G_{ij}>0]$ captures discrete threshold effects at the status boundary independently of gap magnitude. We include $\kappa_g$ to capture a possible discrete effect at the status boundary itself---crossing above the current holder may matter beyond the size of the gap---identified by the jump in the hazard at parity ($G_{ij}=0$), net of the continuous slopes. All channel terms are fully interacted with $g$, yielding slope and discontinuity parameters specific to each skill class.
 
Equation~\ref{eq:gravity_main} states the conceptual full specification; in estimation, its source and target fixed effects cannot enter together. Because the status gap is a linear function of source and target occupational status, source and target fixed effects cannot be absorbed simultaneously \parencite{FafchampsGubert2007, HeadMayer2014}. We therefore estimate the model under two complementary fixed-effect strategies, each identifying the status effect from a different source of variation \parencite{BeineBertoliFernandezHuertas2016}. The first retains source and skill fixed effects $(\alpha_i, \alpha_s)$, identifying the gap effect from within-source-skill, between-target variation and conditioning on all stable characteristics of the source occupation and the skill---so that the estimated effect reflects how target status, relative to a fixed source, governs adoption or abandonment. The second retains target and skill fixed effects $(\alpha_j, \alpha_s)$, identifying the gap effect from within-target-skill, between-source variation and conditioning on all stable characteristics of the target occupation and the skill---so that the estimated effect reflects how source status, relative to a fixed target, shapes the same outcomes. Because each strategy leaves one endpoint uncontrolled, agreement in the sign of estimates across the two panels provides suggestive evidence against endpoint-specific confounding; magnitudes are more endpoint-dependent, with several coefficients attenuating substantially under Panel B and one term (the downward-gap slope for sensory-physical abandonment) differing in sign, which we address directly in the SI. Full coefficient tables for both flows under both specifications are reported in \hyperref[sec:SI_coef_adopt]{Supplementary Sections~S2.2.1--S2.2.2}.

\subsection*{Inference and Robustness}

Standard errors are clustered three-way by source occupation, target occupation, and skill to account for network dependence in the triadic data structure. All main models are estimated on the full risk sets ($n \approx 21.5$M adoption dyads; $n \approx 18.6$M abandonment dyads). Our results are robust to a battery of specification checks, each targeting a distinct threat to inference. RCA denominator-drift checks address the fact that RCA is a relative index whose occupation-weighted mean equals one by construction, which could in principle generate the directional pattern mechanically rather than through occupational behavior: re-estimating the full model with the economy-wide RCA denominator frozen at its 2015 value, and separately with the raw, unnormalized change in O*NET importance ratings as the outcome, preserves the sign of nearly all estimated coefficients---the lone exception, a near-zero, non-significant downward-gap slope for sensory-physical abandonment, is addressed in the SI---with only modest magnitude attenuation among sign-preserved coefficients under the frozen-denominator specification (\hyperref[sec:SI_rca_denom]{Supplementary Section~\ref*{sec:SI_rca_denom}}). A within-stratum permutation test targets the leading alternative---that a common technological shock, not directional diffusion, drives the pattern: permuting outcomes within skill-type $\times$ source-status-quintile cells breaks the link between the status gap and skill change while holding stratum rates fixed, and the observed coefficients fall far outside the resulting null distribution ($B = 1{,}000$ permutations)---showing the directional signal is not an artifact of stratum-level rates and is hard to reconcile with a common shock, which would not on its own reverse across skill domains (\hyperref[sec:SI_permutation]{Supplementary Section~\ref*{sec:SI_permutation}}). RCA threshold sensitivity analyses reconstruct adoption and abandonment risk sets at cutpoints $\{0.90, 1.00, 1.10, 1.25\}$ to confirm that findings do not depend on the exact binarization of skill specialization (\hyperref[sec:SI_rca_threshold]{Supplementary Section~\ref*{sec:SI_rca_threshold}}). Alternative status measure checks re-estimate the gravity model substituting PC1 with each of its three component variables individually---log wage, log educational requirements, and cognitive task content---confirming that the directional asymmetry is invariant to the specific operationalization of occupational status (\hyperref[sec:SI_alt_status]{Supplementary Section~\ref*{sec:SI_alt_status}}). Skill-archetype misclassification checks reassign a random fraction $p \in \{0.10, 0.20\}$ of skills to a new archetype drawn from the empirical marginal distribution across $B = 200$ replications, confirming that the directional signs are overwhelmingly preserved across all coefficient $\times$ panel combinations (\hyperref[sec:SI_missclass]{Supplementary Section~\ref*{sec:SI_missclass}}). Sub-period stability checks replicate Panel~A estimates across the 2015--2018, 2019--2021, and 2022--2024 windows to verify that directional asymmetries are not driven by a single phase of the observation window (\hyperref[sec:SI_subperiod]{Supplementary Section~\ref*{sec:SI_subperiod}}). Robustness checks replacing cosine distance with skill-network shortest-path and resistance distances leave conclusions unchanged. All estimation uses \texttt{feglm} from the \texttt{fixest} package in R \parencite{berge_fixest_2025}.

\subsection*{Use of AI-assisted technologies}

We used Anthropic's Claude in two ways: Claude (Opus~4.8, accessed via Claude Code) to assist with analysis-code development and correctness checking and to audit the manuscript's internal consistency against the replication code, and Claude (Sonnet~4.6, accessed via the Claude chat interface) to improve the grammar, clarity, and concision of author-written text. These tools generated no data, results, or scientific claims; all analyses were specified, executed, and interpreted by the authors, who reviewed all AI-assisted output for accuracy and potential bias and are solely responsible for the final content. No AI tool is listed as an author or cited as a source.


\section*{Acknowledgments}

\textbf{General:} Our use of AI-assisted technologies is described above (see Methods); the authors are solely responsible for the content of this work.

%\textbf{Funding:} ANID Chile, FONDECYT Regular grant 1261889 (M.B.).

%\textbf{Author contributions:} Conceptualization: R.C. and M.B. Formal analysis: R.C. and M.B. Funding acquisition: M.B. Methodology: R.C. and M.B. Writing---original draft: R.C. and M.B. Writing---review and editing: R.C. and M.B.

%\textbf{Competing interests:} The authors declare that they have no competing interests.

\textbf{Data and materials availability:} All data needed to evaluate the conclusions of the paper are present in the paper and/or the Supplementary Materials. The analysis draws on publicly available sources: the O*NET database (2015--2024; \url{https://www.onetcenter.org/database.html}), the Bureau of Labor Statistics Occupational Employment and Wage Statistics (\url{https://www.bls.gov/oes/}), and the O*NET--SOC crosswalk files. All derived datasets and the complete analysis code---data processing, model estimation, and figure and table generation are deposited at \href{https://anonymous.4open.science/r/skill_diffusion_replication-7857/README.md}{\nolinkurl{replication_package.html}}. No restrictions apply to the availability of materials.

\textbf{Interactive supplementary tool.} An interactive explorer of the individual directed triads $(i,j,s)$ underlying the gravity model estimates---organized by log-hazard tercile and status-gap zone---is available at \href{https://anonymous.4open.science/w/skill_diffusion_replication-7857/triads_explorer.html}{\nolinkurl{triads_explorer.html}}. (see \hyperref[sec:SI_triads]{Supplementary Section~\ref*{sec:SI_triads}}).

\clearpage
\bibliographystyle{naturemag}
\bibliography{references}

%%%%%%%%%%%%%%%%%%%%%%%%%%%%%%%%%%%%%%%%%%%%%%%%%%%%%%%%%%%%%%%%%%%%%%%%%%%%%%%%%
\clearpage
\appendix

\renewcommand{\thesection}{S\arabic{section}}
\renewcommand{\thesubsection}{S\arabic{section}.\arabic{subsection}}
\renewcommand{\thesubsubsection}{S\arabic{section}.\arabic{subsection}.\arabic{subsubsection}}
\renewcommand{\thetable}{SM\arabic{table}}
\renewcommand{\thefigure}{SM\arabic{figure}}
\setcounter{section}{0}
\setcounter{table}{0}
\setcounter{figure}{0}

\addtocontents{toc}{\protect\setcounter{tocdepth}{3}}
\renewcommand{\contentsname}{Supplementary Materials --- Table of Contents}
\tableofcontents
\vspace{1cm}

\clearpage

\begin{center}
    {\large\textbf{Supplementary Materials}}\\[0.5em]
    {\normalsize "A Matthew Effect in Occupational Skill Content: \\
    Divergent Diffusion of Cognitive and Physical Skills Reinforces Inequality"}
\end{center}
\vspace{1em}

\section{Methods}

\subsection{Data and Measures}

\subsubsection{O*NET Extracts, Rolling Update Design, and Taxonomic Harmonization}
\label{sec:SI_onet}
 
We use the O*NET Database, a nationally administered survey of occupational skill requirements maintained by the U.S.\ Department of Labor. O*NET reports Importance and Level ratings for 160 skill, knowledge, and ability descriptors for each of approximately 873 standardized occupations via a rolling update design in which approximately 20\% of occupations are resurveyed each year. Each occupation is assigned to an annual cohort; within a cohort, data collection proceeds over 18--24 months through employer surveys. This design means that, in any given release, a majority of occupations carry forward scores from a prior release rather than reflecting current observation. We use the Importance dimension, which measures the degree to which a given requirement is needed to perform the occupation's core tasks; this choice follows \textcite{alabdulkareem_unpacking_2018} and avoids conflating routine task difficulty (Level) with the need for the requirement.
 
For this study we define baseline ($t_0 = 2015$) as the O*NET release in which an occupation was most recently actively updated within calendar year 2015, and endline ($t_1 = 2024$) as the most recent O*NET release available as of December 2024 in which an occupation was actively updated on or after January 2022, ensuring at least one full refresh cycle relative to baseline. We classify an occupation as \emph{actively updated} in year $y$ if its update date in the O*NET release falls within calendar year $y$. This endpoint-window design is preferred over an annual panel for three reasons. First, intermediate-wave non-change is ambiguous---it may reflect genuine stability or measurement latency; an annual-panel design would interpret latency as zero adoption, producing false negatives. Second, the nine-year window guarantees that all occupations receive at least one full survey refresh, allowing adoption to be measured cleanly at endline. Third, the endpoint design is consistent with the interval-censored hazard framework: we observe whether adoption occurred within $(t_0, t_1]$ without imputing annual timing.
 
Between 2015 and 2024 the O*NET taxonomy underwent several revisions: SOC codes were split, merged, or discontinued. We resolve these using the O*NET--SOC crosswalk files. Direct matches are retained as-is. For splits (one 2015 code $\to$ $\geq$2 successor codes), 2015 scores are assigned to each successor and adoption in any successor counts; results are robust to alternative split-handling rules, including employment-weighted aggregation and requiring adoption in the modal successor. For merges ($\geq$2 codes $\to$ one), the predecessor with the highest 2015 importance rating is retained for each descriptor. Discontinued codes are excluded from the risk set. After harmonization, the raw taxonomy covers 873 occupations present at both endpoints; the analytic sample of 741 occupations is derived from this set by additional restrictions documented in \hyperref[sec:SI_sample]{Supplementary Section~\ref*{sec:SI_sample}}.

\subsubsection{Wages, Education, and Task-Based Socio-Cognitive Intensities}
\label{sec:SI_wages}
 
Occupational wages are drawn from the Bureau of Labor Statistics Occupational Employment and Wage Statistics (OEWS) for 2015. We use the annual median wage (\texttt{A\_MEDIAN}) at the detailed occupational level, log-transformed to reduce skew. Educational requirements are operationalized using the O*NET \emph{Required Level of Education} scale (Scale ID: RL), which asks incumbent workers to report the level of education typically required for entry into their occupation across twelve ordered categories ranging from less than a high school diploma to post-doctoral training. We compute the occupation-level expected required-education level as $\mathrm{edu}_j = \sum_{c=1}^{12} c\,p_{jc} \big/ \sum_{c=1}^{12} p_{jc}$, where $c$ indexes the twelve ordered categories and $p_{jc}$ is the share of surveyed workers in occupation $j$ reporting category $c$; the resulting score (range $[1,12]$) is log-transformed before entry into the status index. Cognitive task content is constructed following \textcite{alabdulkareem_unpacking_2018}: for each occupation $j$ we compute the fraction of total O*NET importance weight accounted for by socio-cognitive skill items,
%
\begin{equation*}
\mathrm{cog}_j = \frac{\displaystyle\sum_{s \,\in\, \text{Cognitive}} \mathrm{onet}(j,s)}{\displaystyle\sum_{s \,\in\, S} \mathrm{onet}(j,s)},
\end{equation*}
%
where $\mathrm{onet}(j,s) = ({\rm IM}_{js} - 1)/4 \in [0,1]$ maps the raw importance rating to the unit interval and $\text{Cognitive}$ is the set of socio-cognitive items identified by the Louvain domain classification described in \hyperref[sec:SI_domain]{Supplementary Section~\ref*{sec:SI_domain}}. The resulting score $\mathrm{cog}_j \in [0,1]$ measures the share of an occupation's total skill requirement that is socio-cognitive in nature, independently of wage and educational attainment. All three variables---log wage, log education, and cognitive task content---are constructed from 2015 values and held fixed throughout the observation window, ensuring they are predetermined with respect to the adoption and abandonment outcomes measured at endline.
 
 \subsection{Skill Classification by Functional Domain and Specificity}
\label{sec:SI_skill_class}
 
The skill classification procedure combines two established approaches. From Alabdulkareem et al.~\parencite{alabdulkareem_unpacking_2018} we adopt the revealed comparative advantage (RCA) threshold for binarizing the occupation--skill matrix and the Louvain community detection algorithm for partitioning skills into functional domains. From Hosseinioun et al.~\parencite{hosseinioun_skill_2025} we adopt the nestedness contribution score $c_s$ for characterizing each skill's structural position within that matrix in terms of directed conditional dependencies. A defining feature of our implementation is that all four analytical components---domain classification via Louvain, the adoption risk set, the abandonment risk set, and the nestedness contribution score---operate on the identical 
binary RCA$\,{\geq}\,1$ matrix. Hosseinioun et al.~\parencite{hosseinioun_skill_2025} binarize their occupation--skill matrix using a disparity filter applied to continuous skill level scores; we use RCA$\,{\geq}\,1$ throughout because it is precisely the threshold that defines skill specialization in the adoption and abandonment outcomes---ensuring that the classification of skills and the construction of risk sets and outcomes rest on a single, coherent criterion. The subsections below describe each component in turn.
 
\subsubsection{Revealed Comparative Advantage (RCA)}
\label{sec:SI_rca}
 
For each occupation--skill pair $(j, s)$, the RCA score measures whether skill $s$ is more intensively required by occupation $j$ than the economy-wide average would predict:
%
\begin{equation*}
\mathrm{RCA}(j, s) = \frac{v_{js} / \sum_{s'} v_{js'}}{\sum_{j'} v_{j's} / \sum_{j'}\sum_{s'} v_{j's'}},
\end{equation*}
%
where $v_{js}$ is the mean O*NET importance score for skill $s$ in occupation $j$, averaged across survey respondents. An occupation is considered specialized in skill $s$ if $\mathrm{RCA}(j,s) \geq 1$ (implemented in the replication code as the strict inequality $\mathrm{RCA}(j,s) > 1$; because exact equality to one does not arise for the continuous importance ratios used here, the two conventions yield identical specialization indicators). This normalization removes skill ubiquity as a confound: a skill required at low intensity by virtually all occupations does not register as a specialization in any of them, whereas a skill disproportionately concentrated in a subset does \parencite{alabdulkareem_unpacking_2018}. All RCA scores and the resulting binary specialization indicators are computed from 2015 O*NET data and held fixed throughout the analysis.
 
\subsubsection{Functional Domain}
\label{sec:SI_domain}
 
We classify each of the 160 skill requirements into one of two functional domains---socio-cognitive or sensory-physical---by applying community detection to the skill co-specialization network, following \textcite{alabdulkareem_unpacking_2018}. We compute pairwise cosine similarities between the binary RCA$\,{\geq}\,1$ specialization vectors of all skill pairs across the 741 occupations, construct an undirected weighted network in which nodes are skills and edge weights are cosine similarities, and retain only edges with weight exceeding 0.1 to remove noise from weakly co-occurring pairs. Applying the Louvain community detection algorithm \parencite{blondel_fast_2008} to this network recovers a stable two-community solution: one cluster concentrates social, cognitive, and knowledge-intensive items (97 skills) and the other concentrates physical, sensory, and manual items (63 skills). We label these communities \textit{socio-cognitive} and \textit{sensory-physical} respectively. Applying the same procedure to the 2024 endline data recovers the identical two-community partition, confirming that the domain structure is stable across the observation window; we use the 2015 classification throughout to ensure domain assignment is predetermined with respect to outcomes.
 
\subsubsection{Functional Specificity}
\label{sec:SI_specificity}
 
Within the binary RCA$\,{\geq}\,1$ matrix established above, we ask a second-order question: not which skills are concentrated in which occupations, but how each skill's position within the matrix organizes conditional dependencies across the full skill space. We quantify this using the nestedness contribution score $c_s$, following \textcite{hosseinioun_skill_2025}. The occupation--skill matrix is nested when the specialization profiles of occupations with narrower portfolios are systematically contained within those of occupations with broader ones---narrowly specialized occupations require only skills that broadly specialized ones also require, but not vice versa. This pattern reflects directed conditional dependencies: the presence of a narrowly applicable skill in an occupational profile presupposes the presence of more broadly applicable ones, but the reverse does not hold. A skill with high $c_s$ organizes these directed dependencies---it is a structural prerequisite that other skills build upon, and its absence forecloses them. A skill with low or negative $c_s$ appears more independently of others and contributes little to the hierarchical dependency structure. We quantify nestedness using the NODF index (Nestedness based on Overlap and Decreasing Fill) \parencite{almeida_neto_consistent_2008} and define the nestedness contribution of skill $s$ as the standardized change in NODF when $s$ is removed from the baseline matrix:
%
\begin{equation*}
c_s = \frac{\mathrm{NODF}_{\text{full}} - \mathrm{NODF}_{-s}}{\mathrm{SD}(\mathrm{NODF}_{\text{permuted}})},
\end{equation*}
%
where $\mathrm{NODF}_{-s}$ is the index after removing skill $s$ and the denominator is the standard deviation of NODF across 5{,}000 randomly permuted matrices preserving row and column marginals, following the procedure of \textcite{hosseinioun_skill_2025}. All $c_s$ scores are computed from the 2015 baseline matrix and are therefore predetermined with respect to outcomes.
 
\subsubsection{Full Three-Class Taxonomy}
\label{sec:SI_taxonomy}
 
Crossing functional domain with functional specificity yields the three-class taxonomy used throughout the analysis. Within the socio-cognitive domain, all 97 skills have $c_s > 0$ in our data---the Louvain community that defines the cognitive cluster coincides with the nested core of the occupation--skill matrix, so that every cognitive skill contributes positively to the hierarchical dependency structure. We split these at the within-domain median of $c_s$: skills below the median are classified as \textit{specialized socio-cognitive} (48 skills)---narrowly concentrated cognitive skills that build on a broader foundation---and skills at or above the median as \textit{general socio-cognitive} (49 skills)---broadly scaffolding skills that underpin the nested structure. Within the sensory-physical domain, all 63 skills are classified as \textit{sensory-physical} regardless of their $c_s$ value: the high-nestedness sensory-physical cell is empirically absent in both 2015 and 2024 data, and results are unchanged when it is retained as a separate fourth category. The distribution of $c_s$ by domain, with the within-domain median cutoff for the socio-cognitive split, is shown in Fig.~\ref{fig:SI_cs_distribution}.

\begin{figure*}[tp]
\centering
\includegraphics[width=0.65\textwidth]{figures/fig_S2_cs_distribution.pdf}
\caption{\textbf{Distribution of nestedness contribution scores ($c_s$) by functional domain.} Each curve shows the kernel density of $c_s$ across skills within one functional domain. Within the socio-cognitive domain (teal), all skills have $c_s > 0$ in our data; the dashed vertical line marks the within-domain median, which defines the boundary between specialized socio-cognitive (left of median, $n = 48$) and general socio-cognitive (right of median, $n = 49$) skills. Within the sensory-physical domain (gray), $c_s$ values are negative throughout, confirming the absence of a high-nestedness sensory-physical cluster. O*NET 2015.}
\label{fig:SI_cs_distribution}
\end{figure*}


 
\subsection{Occupational Status Index}
\label{sec:SI_status}
 
\subsubsection{Component Variables, Scaling, and PCA Construction}
\label{sec:SI_pca}
 
We construct a scalar status index $\sigma_i$ for each of the 741 occupations in the analytic sample by applying principal component analysis (PCA) to three occupation-level variables: median annual wage (log-transformed), the expected level of required education (O*NET Scale~RL, log-transformed), and cognitive task content. The first two variables are drawn from the Bureau of Labor Statistics OEWS and O*NET respectively, as described in \hyperref[sec:SI_wages]{Supplementary Section~\ref*{sec:SI_wages}}. Cognitive task content is the fraction of total O*NET importance weight accounted for by socio-cognitive skill items, as defined in \hyperref[sec:SI_wages]{Supplementary Section~\ref*{sec:SI_wages}}. All three variables are measured at baseline (2015) and held fixed throughout the observation window.
 
We estimate the PCA with mean-centering and unit-variance scaling (\texttt{center = TRUE}, \texttt{scale. = TRUE}) across all 741 occupations. This standard specification ensures that each variable contributes to PC1 in proportion to its covariation with the others, rather than in proportion to its raw variance. The first principal component (PC1) captures the large majority of the joint variance across the three dimensions, with loadings of comparable magnitude across all three variables. The near-equal loading structure confirms that wages, educational requirements, and cognitive task content converge on a single underlying status axis rather than one variable dominating the composite. We orient PC1 so that higher values correspond to higher status---verified by confirming that dyads in which the target outranks the source on all three raw variables yield a positive gap, peer dyads yield a gap near zero, and dyads in which the target ranks below the source on all three yield a negative gap (Fig.~\ref{fig:SI_pca_status}).
 
The signed status gap for each directed dyad is then computed as the difference in occupation-level scores: $G_{ij} = \sigma_j - \sigma_i$, where $i$ is the source occupation and $j$ is the target. Positive values of $G_{ij}$ indicate that the target ranks above the source; negative values indicate the reverse. The full distribution of occupation-level status scores, the PCA biplot, and pairwise correlations among the three component variables are reported in \texttt{output/tables/main/pca\_status\_decision.csv} and the accompanying figure \texttt{output/figures/main/fig\_pca\_status.pdf}, both generated by \texttt{04\_status\_pca.R} in the replication repository.

\begin{figure*}[tp]
\centering
\includegraphics[width=0.95\textwidth]{figures/fig_pca_status.pdf}
\caption{\textbf{Principal component analysis of the occupational status index.}
\textbf{(A)} PCA loadings for the three principal components extracted from log-transformed median wage (BLS OEWS 2015), log-transformed expected level of required education (O*NET Scale RL), and cognitive task content (fraction of O*NET importance weight accounted for by socio-cognitive skill items), estimated with mean-centering and unit-variance scaling across all 741 occupations in the analytic sample. PC1 captures 75.0\% of joint variance, with loadings of comparable magnitude across all three variables (log wage: $|{-}0.554|$; log education: $|{-}0.629|$; cognitive score: $|{-}0.545|$), confirming convergence on a single underlying status axis. The sign of PC1 is reversed (flip applied) so that higher values correspond to higher occupational status. \textbf{(B)} Distribution of occupation-level status scores ($\sigma_i$) across the 741 analytic occupations. The dashed vertical line marks the origin ($\sigma_i = 0$), which corresponds to the sample mean after mean-centering. The signed status gap for each directed dyad is computed as $G_{ij} = \sigma_j - \sigma_i$.}
\label{fig:SI_pca_status}
\end{figure*}


\subsection{Analytic Sample}
\label{sec:SI_sample}
 
The analytic sample comprises 741 occupations derived from the 873 standardized O*NET occupations present in both the 2015 baseline and 2024 endline extracts after taxonomic harmonization (\hyperref[sec:SI_onet]{Supplementary Section~\ref*{sec:SI_onet}}). The single restriction applied is that each occupation must have an actively updated O*NET profile at both endpoints: it must have received a full survey refresh within the 2015 baseline cohort and at least one additional refresh between January 2022 and December 2024. Occupations present in only one endpoint are excluded from both risk sets, as adoption and abandonment cannot be measured cleanly without valid baseline and endline profiles. This restriction is implemented as the intersection of the sets of occupations active in both O*NET extracts (\texttt{intersect(occ\_2015, occ\_2024)} in \texttt{01a\_risk\_set\_adoption.R} and \texttt{01b\_risk\_set\_abandonment.R}), yielding 741 occupations.
 
All 741 occupations have complete covariate coverage on all three status-index components---median wage, educational requirements, and cognitive task content---so that the occupational status index $\sigma_i$ is defined without exception across the full analytic sample. At the dyadic level, wage and cognitive gaps show partial coverage (88.9\% and 89.2\% of directed dyads, respectively) owing to imperfect matching between BLS OEWS occupational codes and the O*NET SOC taxonomy; educational gaps show near-complete coverage (99.3\%). These dyadic missingness patterns affect the gap variables used as diagnostic covariates but do not affect the status index, which is constructed from occupation-level scores derived independently of the dyadic crosswalk. All 741 occupations appear as both sources and targets in the adoption and abandonment risk sets, confirming that no systematic truncation of the occupational hierarchy occurs at either end of the status distribution.
 
\section{Findings}

\subsection{Descriptive Statistics}
\label{sec:SI_descriptives}
 
Table~\ref{tab:SI_occ_desc} summarizes the occupation-level variables used to construct the status index across the 741 occupations in the analytic sample. The sample spans the full range of the U.S. occupational structure: from low-wage food-preparation work to high-wage physicians and surgeons, and from occupations requiring less than a high school diploma to those requiring post-doctoral training. The status index (PC1) has mean zero by construction and exhibits substantial dispersion in occupational standing across the hierarchy.
 
Table~\ref{tab:SI_riskset_desc} describes the adoption and abandonment risk sets by skill class. The adoption risk set comprises approximately 21.5 million directed dyads across 160 skill requirements; the abandonment risk set comprises approximately 18.6 million dyads. Baseline adoption rates are highest for general socio-cognitive skills and lowest for sensory-physical skills, reflecting the greater mobility of cognitive skills across occupational boundaries. Baseline abandonment rates are uniformly higher than adoption rates across all skill types, indicating that skill loss is more common than skill gain over the 2015--2024 window — a pattern consistent with the overall restructuring of occupational content documented in the main text.
 
\input{tables/tab_SI_occ_desc.tex}
 
\input{tables/tab_SI_riskset_desc.tex}


\subsection{Fixed-Effects Gravity Models: Status Gaps Drive Diffusion, Net of Occupational and Skill Heterogeneity}
\label{sec:SI_coef_tables}

Tables~\ref{tab:SI_coef_adopt} and~\ref{tab:SI_coef_aband} report complete coefficient estimates from the complementary log-log gravity models for adoption and abandonment under both fixed-effect strategies. Panel~A retains source and skill fixed effects $(\alpha_i, \alpha_s)$, identifying the status gap effect from within-source-skill variation across targets; Panel~B retains target and skill fixed effects $(\alpha_j, \alpha_s)$, identifying the effect from within-target-skill variation across sources. Both panels condition on cosine profile distance, so that the directional gap terms index movement along the status hierarchy net of task similarity. The status boundary effect captures a discrete shift in the log-hazard at $G_{ij} = 0$, independently of gap magnitude. All channel terms are fully interacted with skill class, yielding separate slope and discontinuity parameters for each of the three skill types. Standard errors are clustered three-way by source occupation, target occupation, and skill to account for network dependence in the triadic data structure.

\subsubsection{Full Coefficient Tables: Adoption}
\label{sec:SI_coef_adopt}

\input{tables/tab_SI_coef_adopt.tex}

\subsubsection{Full Coefficient Tables: Abandonment}
\label{sec:SI_coef_aband}

\input{tables/tab_SI_coef_aband.tex}

One coefficient in Table~\ref{tab:SI_coef_aband} differs in both sign and significance between panels: the downward-gap slope for sensory-physical abandonment ($\hat{\beta}^{\downarrow}$) is small and non-significant under Panel~A but positive and significant under Panel~B. We do not interpret this specific term as robust. Our claims about sensory-physical abandonment in the main text rest instead on the upward-gap slope ($\hat{\beta}^{\uparrow}$), which is positive, significant, and consistently signed under both specifications.


\subsection{Skill-Level Divergent Diffusion Reproduces Aggregate Occupational Stratification}
\label{sec:SI_projections}
 
Figure~\ref{fig:projection_main} of the main text reports the projection under both fixed-effect strategies (Panel~A: source and skill fixed effects; Panel~B: target and skill fixed effects), each against two nested counterfactual nulls that decompose the sources of the observed gradient. This section details those nulls; the full Q5--Q1 gradient statistics are reported in Table~\ref{tab:SI_gradients}.
 
The two nulls are nested within the full directional model, each adding one component of the directional diffusion argument. The distance-only null applies only the skill profile distance parameter ($\hat{\delta}$), setting all status gap terms to zero; it asks whether task similarity alone can account for the stratified pattern of requirement change. The symmetric status null applies the skill profile distance parameter and a single symmetric coefficient on the absolute value of the status gap ($b_{\mathrm{avg}} \times |G_{ij}|$), asking whether the \textit{size} of the status gap---irrespective of its direction---accounts for the gradient beyond proximity. The full directional model applies all parameters ($\hat{\beta}^{\uparrow}$, $\hat{\beta}^{\downarrow}$, $\hat{\kappa}$, $\hat{\delta}$), capturing the asymmetric effect of gap direction.
 
The results reveal a clear decomposition. The distance-only null stays far below the directional model across all skill types and both flows, recovering only a small fraction of each observed gradient and never approaching the model's recovery. The symmetric status null performs only marginally better, confirming that the size of the status gap without direction carries little explanatory power---particularly for abandonment, where it recovers almost none of the observed socio-cognitive gradient under Panel~A. Table~\ref{tab:SI_gradients} reports each benchmark's gradient (columns Model, Sym, Null) alongside the model's advantage expressed as a ratio (final two columns). We read that advantage as the absolute gap between the directional model and the symmetric null---directly comparable across the Model and Sym columns---rather than as the ratio, because the ratio is unstable: where the symmetric null's gradient is near zero, small absolute changes translate into multipliers that swing wildly across fixed-effect strategies (reaching $50\times$--$131\times$ for socio-cognitive abandonment under Panel~A), whereas the absolute gap between the directional model and the symmetric null is modest but stable across specifications. The full directional model recovers most of each observed Q5--Q1 gradient under source fixed effects (Panel~A) and a smaller share under target fixed effects (Panel~B), as the gap coefficients attenuate once the target's own status is absorbed. These results confirm that the gradient of occupational skill consolidation is driven not by proximity or by the magnitude of status differences, but by their direction; the magnitude of that advantage over the symmetric null is best read in absolute, not ratio, terms. The full quantitative breakdown---observed, model, and both null gradients with recovery and advantage figures---is reported in Table~\ref{tab:SI_gradients}.
 
\input{tables/tab_SI_gradients.tex}
 
% Both projection panels (source-FE Panel A and target-FE Panel B) now appear
% combined under a common caption as Fig.~\ref{fig:projection_main} in the main text.
 

\subsection{Predicted triads by diffusion direction and skill type}
\label{sec:SI_triads}

We provide an interactive explorer that renders the individual directed triads $(i, j, s)$ underlying the directional gravity model, available at \href{https://anonymous.4open.science/w/skill_diffusion_replication-7857/triads_explorer.html}{\nolinkurl{triads_explorer.html}}.
The explorer organizes triads into a $3 \times 3$ matrix whose rows correspond to log-hazard terciles (High / Mid / Low) and whose columns correspond to status-gap zones: T1 (target ranks below source, $G_{ij} < 0$), T2 (peer / lateral, $G_{ij} \approx 0$), and T3 (target ranks above source, $G_{ij} > 0$). Within each cell, triads are ranked by their log-partial-hazard $\ell p_{ij} = \kappa_{g}\,\mathbbm{1}[G_{ij}>0] + \beta^{\uparrow}_{g}\Delta^{\uparrow}_{ij} + \beta^{\downarrow}_{g}\Delta^{\downarrow}_{ij}$, where higher values indicate greater predicted probability of adoption or abandonment. Each entry reports the modal source--target occupational pair for that skill among realized diffusion events in O*NET 2015--2024. Controls allow selection of flow (adoption or abandonment), fixed-effect strategy (Source $+$ Skill FE or Target $+$ Skill FE), structural distance range, sampling method, and skill type. A scatter plot shows each skill's mean log-hazard against its mean status gap across realized events, with point size proportional to triad frequency.

\section{Robustness Checks}

\subsection{Robustness to RCA's Relative Construction}
\label{sec:SI_rca_denom}

Revealed comparative advantage is a relative (Balassa-style) index: for any skill, the occupation-mass-weighted average of RCA across all occupations in the analytic sample equals exactly one by construction. This means that if real adoption of a skill rises among high-status occupations, the economy-wide denominator against which every other occupation's specialization is measured necessarily rises as well, which could in principle push low-status occupations below the $\mathrm{RCA} \geq 1$ threshold without any change in their own underlying behavior---generating spurious abandonment events among low-status occupations and, by symmetry, spurious adoption events at the upper end, as an artifact of the index's construction rather than of directional diffusion. We test this directly by holding the economy-wide denominator fixed at its 2015 value while letting each occupation's own specialization share update normally through 2024, isolating real behavioral change from denominator drift. We additionally estimate an alternative specification that discards RCA normalization entirely and uses the raw, unnormalized change in O*NET importance ratings as the dependent variable, z-scored within each flow prior to estimation so that coefficients are expressed in standard-deviation units and directionally comparable across specifications. Because adoption and abandonment are modeled as separate directional events under the RCA framework whereas raw importance change is a single continuous variable in which abandonment corresponds to a decrease, the sign of the raw-importance coefficients for the abandonment flow is multiplied by $-1$ before comparing directions. All remaining model components, including the full set of structural fixed effects and skill-class interactions in Equation~\ref{eq:gravity_main}, are left unchanged across all three specifications.

Results are reported for both Panel~A (source~$+$~skill fixed effects) and Panel~B (target~$+$~skill fixed effects) in Table~\ref{tab:SI_rca_denom}. Freezing the denominator preserves the sign of 11 of 12 estimated slope parameters in Panel~A and all 12 in Panel~B, with coefficient magnitudes remaining stable across specifications. The single Panel~A exception is the downward-gap abandonment slope for sensory-physical skills ($\hat{\beta}^{\downarrow} = -0.006$, SE~$= 0.040$), which is statistically indistinguishable from zero in the standard specification and does not constitute a meaningful sign reversal. Notably, for the upward-gap abandonment slope of sensory-physical skills, freezing the denominator sharpens the estimate from $0.095$ to $0.107$ in Panel~A, confirming that mechanical redistribution of the baseline index modestly understates the rate at which higher-status occupations shed manual content. The z-scored raw-importance specification preserves the sign of 8 of 12 slope parameters in Panel~A and 11 of 12 in Panel~B. The four Panel~A discrepancies follow a consistent pattern: the downward-gap slopes for both socio-cognitive archetypes in adoption reverse sign, and two near-zero coefficients in the abandonment flow produce nominal sign disagreements that are not statistically meaningful. The single Panel~B discrepancy involves the General~SC downward-gap adoption slope, whose standard estimate is $-0.001$ and is indistinguishable from zero. The core results---the positive $\hat{\beta}^{\uparrow}$ for both socio-cognitive archetypes in adoption and the full sensory-physical pattern across both flows and panels---are reproduced under the raw-importance specification throughout. Together, these results indicate that the status-asymmetric gradient documented in the main text reflects real occupational behavior rather than the redistribution properties of a relative index.

\input{tables/tab_SI_rca_denom.tex}

\subsection{Robustness to Implicit Source-Multiplicity Weighting}
\label{sec:SI_weighting}

The standard riskset contains one row per directed dyad $(i, j, s)$, where $i$ is a baseline holder of skill $s$ and $j$ is the focal occupation. A realized event $(j, s)$ therefore contributes $n_{js}$ rows to the likelihood, where $n_{js} = |S_s|$ is the number of eligible sources for that target--skill pair. The skill fixed effect $\alpha_s$ absorbs the average source count per skill but not the within-skill, between-target variation in $n_{js}$, which is precisely the variation from which the Panel~A gap slopes $\hat{\beta}^{\uparrow}$ and $\hat{\beta}^{\downarrow}$ are identified. If source-rich targets sit systematically at one end of the status hierarchy, the estimated coefficients could in principle reflect row multiplicity rather than occupational behaviour.

We address this directly by re-estimating the main model with frequency weights $w_{ij} = 1/n_{js}$, so that the total contribution of each realized event $(j, s)$ to the likelihood equals
$\sum_{i \in S_s} 1/n_{js} = 1$ regardless of how many sources it has. The formula, fixed-effect structure, and three-way clustering are otherwise identical to Tables~\ref{tab:SI_coef_adopt} and~\ref{tab:SI_coef_aband}. In the analytic sample, the source pool per target--skill cell has a
median of $346$ sources for adoption and $352$ for abandonment (maximum $416$ and $415$, respectively), so the implicit weighting that the unweighted model applies is substantial and the check is non-trivial.

Table~\ref{tab:SI_weighting} reports the weighted estimates alongside the unweighted baseline for all six slope parameters per flow and panel. All 24 slope parameters preserve sign under frequency weighting, and coefficient magnitudes change by at most $11.8\%$ relative to the unweighted specification (mean change $3.7\%$). The largest individual change is the Physical adoption $\hat{\beta}^{\downarrow}$ in Panel~A ($-0.076$ unweighted vs.\ $-0.085$ weighted), which moves in the direction of a stronger downward barrier to physical-skill adoption, reinforcing rather than undermining the main finding. Together, these results confirm that the status-asymmetric gradient is not an artifact of implicit source-multiplicity weighting.

\input{tables/tab_SI_weighting.tex}

\subsection{Within-Stratum Permutation Test}
\label{sec:SI_permutation}

The leading alternative to directional diffusion is a correlated shock: a common technological force that independently pushes many occupations toward the same skills, so that the apparent role of source status is spurious. We test this by permuting the outcome within strata defined by skill type and source-status quintile. This holds each stratum's overall rate of skill change fixed while breaking any systematic link between the signed status gap and which occupations gain or shed a skill---the structure implied if the status gap were irrelevant and a shared shock drove the change. Across $B = 1{,}000$ permutations, the observed Panel~A coefficients lie between 8 and 211 null standard deviations beyond the permuted distribution (all $p < 0.001$; adoption: $\hat{\beta}^{\uparrow}_{\text{Gen. SC}}$ at $+87$ SDs, $\hat{\beta}^{\uparrow}_{\text{Spec. SC}}$ at $+84$, $\hat{\beta}^{\uparrow}_{\text{Phys}}$ at $-211$; abandonment: $-179$, $-166$, and $+62$, respectively). This indicates the directional signal is not an artifact of stratum-level rates of change. It is also hard to attribute to a common shock, which would have to be not merely directional but \emph{domain-reversing}---driving cognitive and physical skills in opposite directions along the same status gap---a structure a shared exogenous force is unlikely to generate on its own (Fig.~\ref{fig:SI_test_iii}).

\begin{figure*}[tp]
\centering
\includegraphics[width=0.95\textwidth]{figures/fig_SI_test_iii.pdf}
\caption{\textbf{Observed coefficients are far outside the within-stratum permutation null.} Each histogram shows the distribution of $B = 1{,}000$ permuted estimates, where the outcome is randomly reassigned within cells defined by skill type $\times$ source status quintile. The vertical line marks the observed Panel~A estimate; the annotation reports the observed value and its distance from the null in null standard deviations. Rows correspond to $\hat{\beta}^{\uparrow}$ (top) and $\hat{\beta}^{\downarrow}$ (bottom); columns to skill type. All twelve observed coefficients fall outside the null distribution ($p < 0.001$), inconsistent with correlated technological shocks as the primary driver of the directional asymmetry. O*NET 2015--2024. \textbf{(A)} Adoption. \textbf{(B)} Abandonment.}
\label{fig:SI_test_iii}
\end{figure*}


\subsection{Measurement Sensitivity}
 
\subsubsection{RCA Threshold Sensitivity}
\label{sec:SI_rca_threshold}

The baseline analysis defines skill specialization via revealed comparative advantage at $\mathrm{RCA} \geq 1.00$. To assess sensitivity to this threshold, we reconstruct the adoption and abandonment risk sets at three alternative cutpoints ($\mathrm{RCA} \geq 0.90$, $1.10$, $1.25$) and compute observed diffusion rates by status gap decile. The directional pattern is preserved across all four thresholds: socio-cognitive skills exhibit higher adoption rates when the target ranks above the source, and lower rates in the reverse configuration, while sensory-physical skills show the mirror gradient. The absolute rates shift with the threshold---higher thresholds select more specialized dyads with correspondingly higher diffusion rates---but the directional asymmetry is present and of consistent sign at every cutpoint tested. Results are displayed in Fig.~\ref{fig:SI_rca}.
 
\begin{figure*}[tp]
\centering
\includegraphics[width=0.95\textwidth]{figures/fig_S1_rca_threshold.pdf}
\caption{\textbf{Directional diffusion gradients are insensitive to the RCA specialization threshold.} Each panel plots the observed diffusion rate as a function of the signed status gap (target $-$ source, PC1) for one of four specialization thresholds ($\mathrm{RCA} \geq 0.90$, $1.00$, $1.10$, $1.25$; rows) and two flows (Adoption, Abandonment; columns). Colors identify the three skill types. Ribbons are 95\% bootstrap confidence intervals ($B = 500$). The directional asymmetry---socio-cognitive rates rising with the gap, sensory-physical rates falling---is present at every threshold in both flows. O*NET 2015--2024.}
\label{fig:SI_rca}
\end{figure*}

 
\subsubsection{Alternative Status Measures}
\label{sec:SI_alt_status}
 
The main analysis operationalizes occupational status as the first principal component of three occupation-level variables---log median wage, log educational requirements, and cognitive task content---estimated with mean-centering and unit-variance scaling (PC1; see \hyperref[sec:SI_pca]{Supplementary Section~\ref*{sec:SI_pca}}). To assess whether the directional asymmetries documented in the main text depend on this composite operationalization, we re-estimate the gravity model substituting PC1 with each of its three components individually. To ensure comparability across measures and with the main specification, each component is standardized to unit variance at the occupation level before computing directional gaps, so that all coefficients are expressed in units of one standard deviation of the status gap.

As shown in Figures \ref{fig:alt_status_adoption} and \ref{fig:alt_status_abandonment}, the core directional gradient observed in the main analysis---a monotone, kinked response to the signed status gap, steeper in one direction than the other---remains remarkably consistent across all three alternative dimensions. The fundamental asymmetry holds: downward status gaps systematically predict higher abandonment hazards and lower adoption probabilities, particularly for socio-cognitive skills. While the structural pattern is invariant, the magnitude of the effects varies depending on the specific dimension of stratification. The cognitive score measure exhibits the steepest gradients---indicating a particularly severe penalty for moving down the cognitive hierarchy---whereas log wage and log education display more moderate, yet structurally identical, slopes. These findings confirm that the directional status penalty is a robust feature of occupational skill diffusion, independent of how vertical stratification is specifically measured.


\begin{figure*}[tp]
    \centering
    \includegraphics[width=0.8\textwidth]{figures/fig_SI_alt_status_Adoption.pdf}
    \caption{\textbf{Alternative Status Measures for Skill Adoption.} Fitted log-hazard of skill adoption as a function of the signed status gap between target and source occupations. Status gaps are computed using three alternative measures: Log wage, Log education, and Cognitive score. To ensure comparability across measures and with the main specification, all variables were standardized ($SD = 1$) at the occupation level prior to calculating the directional gaps. Panel A displays estimates controlling for source and skill fixed effects, whereas Panel B controls for target and skill fixed effects. Lines represent the estimated upward ($\beta^{\uparrow}$) and downward ($\beta^{\downarrow}$) slopes for each skill archetype. Standard errors are clustered three-way (source, target, and skill).}
    \label{fig:alt_status_adoption}
\end{figure*}



\input{figures/fig_SI_alt_status_Abandonment.tex}

\subsubsection{Robustness to Skill-Archetype Misclassification}
\label{sec:SI_missclass}

The three-class taxonomy assigns each of the 160 skill requirements to one of three archetypes---specialized socio-cognitive, general socio-cognitive, or sensory-physical---based on community detection and nestedness scores computed from the 2015 baseline matrix. To assess whether the directional results depend on the precision of this classification, we conduct a misclassification robustness check in which a random fraction $p \in \{0.10, 0.20\}$ of skills is reassigned to a new archetype drawn from the empirical marginal distribution across the three classes ($48/160$, $49/160$, $63/160$), preserving the overall class proportions in expectation. We conduct $B = 200$ replications at each contamination level, re-estimating the full gravity model under both Panel~A and Panel~B for each perturbed taxonomy, and extract $\hat{\beta}^{\uparrow}$ and $\hat{\beta}^{\downarrow}$ for each archetype across replications.

Across all 24 coefficient $\times$ contamination-level $\times$ panel combinations, sign preservation exceeds 93.5\%, with 23 of 24 combinations achieving 100\% sign preservation. The single exception---the upward-gap slope for sensory-physical abandonment under Panel~A at 20\% contamination---retains the correct sign in 93.5\% of replications, above the conventional 90\% threshold. Mean estimates under 10\% contamination attenuate modestly relative to the baseline (e.g., $\hat{\beta}^{\uparrow}_{\text{Phys, adopt}} = -0.300$ vs.\ $-0.322$ at baseline), with attenuation increasing slightly at 20\% contamination as expected. The directional structure is preserved throughout: socio-cognitive upward slopes remain positive and sensory-physical upward slopes remain negative for adoption, with the reverse pattern for abandonment, under both fixed-effect strategies and both contamination levels. These results confirm that the three-class taxonomy is not a fragile classification: random mislabeling of up to 20\% of skills does not alter the directional pattern documented in the main text. Results are displayed in Fig.~\ref{fig:SI_missclass}.

\begin{figure*}[tp]
\centering
\includegraphics[width=0.75\textwidth]{figures/fig_archetype_misclassification.pdf}
\caption{\textbf{Directional gravity estimates are robust to skill-archetype misclassification.} Each violin shows the distribution of $\hat{\beta}^{\uparrow}$ (top row) and $\hat{\beta}^{\downarrow}$ (bottom row) across $B = 200$ replications in which a random fraction $p$ of skills is reassigned to a new archetype drawn from the empirical marginal distribution. Columns distinguish contamination levels (10\% and 20\% mislabeled); rows within each section distinguish $\hat{\beta}^{\uparrow}$ and $\hat{\beta}^{\downarrow}$. Internal lines in each violin mark the 5th, 50th, and 95th percentiles; the superimposed boxplot shows the interquartile range. Sections correspond to the four flow $\times$ fixed-effect combinations: Adoption Panel~A (blue, top), Adoption Panel~B (blue), Abandonment Panel~A (red), Abandonment Panel~B (red, bottom). Sign preservation exceeds 93.5\% in all 24 combinations and equals 100\% in 23 of 24. O*NET 2015--2024.}
\label{fig:SI_missclass}
\end{figure*}


\subsubsection{Subsample Stability}
 
The robustness checks reported in this section are estimated on a 50\% random subsample of source occupations due to memory constraints, while all main models use the full risk sets. To verify that results do not depend on the specific draw, we re-estimate Panel~A and Panel~B under three independent random seeds (42, 123, 999) for both flows. Across the 36 coefficient $\times$ seed combinations, estimates are effectively flat: the coefficient of variation across seeds is below 5\% in every case, and the absolute standard deviation across seeds is below 0.02 for all coefficients. The directional sign is preserved without exception. Results are displayed in Fig.~\ref{fig:SI_stab}.
 
\begin{figure*}[tp]
\centering
\includegraphics[width=0.95\textwidth]{figures/fig_SI_subsample_stability.pdf}
\caption{\textbf{Directional gravity estimates are stable across independent subsamples.} Each panel plots coefficient estimates from three independent 50\% random draws of source occupations (seeds 42, 123, 999). Rows correspond to $\hat{\beta}^{\uparrow}$ (top) and $\hat{\beta}^{\downarrow}$ (bottom); columns to skill type. Colors distinguish fixed-effect strategies. Near-flat lines indicate that estimates do not depend on the specific subsample drawn. The coefficient of variation across seeds is below 5\% in all cases. O*NET 2015--2024. \textbf{(A)} Adoption. \textbf{(B)} Abandonment.}
\label{fig:SI_stab}
\end{figure*}

 

\subsubsection{Sub-Period Stability of Directional Asymmetries}
\label{sec:SI_subperiod}

The nine-year observation window spans three qualitatively distinct phases: a pre-disruption baseline (2015--2018), a period of pandemic-driven reorganization of work (2019--2021), and the initial diffusion of generative AI (2022--2024). If the observed pattern were driven by correlated technological shocks specific to one of these phases rather than reflecting a structural property of skill diffusion, the directional asymmetry should intensify or reverse in the disruption periods relative to the baseline. Sub-period estimates from Panel~A are consistent with stability rather than shock-driven amplification. For adoption, $\hat{\beta}^{\uparrow}_{\text{Gen. SC}}$ is $+0.310^{***}$ (2015--2018), $-0.019^{\text{ns}}$ (2019--2021), and $+0.300^{***}$ (2022--2024); $\hat{\beta}^{\uparrow}_{\text{Phys}}$ is $-0.189^{***}$, $-0.827^{***}$, and $-0.203^{*}$ across the same windows. The attenuation of $\hat{\beta}^{\uparrow}_{\text{Gen. SC}}$ in 2019--2021 reflects disruption to occupational activity patterns during the pandemic observation period, but the negative sign of $\hat{\beta}^{\uparrow}_{\text{Phys}}$ is preserved and markedly strengthened in that window, confirming that the domain-reversing structure of the gradient is not a feature of any single phase. For abandonment, $\hat{\beta}^{\uparrow}_{\text{Phys}}$ is $+0.188^{**}$, $+0.186^{**}$, and $+0.208^{*}$ across the three windows, with no systematic intensification in the disruption periods, while the socio-cognitive gradient strengthens in 2019--2021 ($\hat{\beta}^{\uparrow}_{\text{Gen. SC}} = -0.434^{***}$, $\hat{\beta}^{\uparrow}_{\text{Spec. SC}} = -0.548^{***}$) before partially attenuating in 2022--2024. Full sub-period estimates are displayed in Fig.~\ref{fig:SI_subperiods}.
 
\begin{figure*}[tp]
\centering
\includegraphics[width=0.95\textwidth]{figures/fig_SI_subperiods.pdf}
\caption{\textbf{Directional friction parameters are stable in direction across sub-periods.} Each panel shows Panel~A estimates ($\hat{\beta}^{\uparrow}$, top row; $\hat{\beta}^{\downarrow}$, bottom row) with 95\% confidence intervals (three-way clustering by source occupation, target occupation, and skill) for three non-overlapping sub-periods. Colors identify the sub-period. The sign of $\hat{\beta}^{\uparrow}$ is preserved across all periods and both flows: positive for socio-cognitive skills in adoption and negative in abandonment; reversed for sensory-physical skills. Wider intervals in 2019--2021 reflect reduced statistical power from the shorter observation window, not reversal of the gradient. The strengthening of $\hat{\beta}^{\uparrow}_{\text{Phys}}$ in adoption and of the socio-cognitive abandonment gradient in 2019--2021 is consistent with pandemic-period acceleration of directional diffusion, but the domain-reversing structure is present in all three windows. O*NET sub-period pairs: 2015--2018 (releases 20.3--23.1), 2019--2021 (24.1--26.1), 2022--2024 (27.1--29.2). \textbf{(A)}~Adoption. \textbf{(B)}~Abandonment.}
\label{fig:SI_subperiods}
\end{figure*}

 

\end{document}
